\section{Compressed additive orbit weights}
\label{sec:weighted-replay}

Orbit counts alone do not describe the individual components of an interval
equivalence relation.  The weighted extension of Agol--Hass--Thurston supplies
the required additional information in compressed form
\cite{AHT2006}.  This section gives a precise implementation contract and
correctness argument for that classical extension.  The contribution here is
its integration with independently checked local-reduction traces, including
signed vector weights and explicit control of representation size.  We do not
claim a new orbit-counting theorem.

\subsection{The weighted output and its representation}
\label{sec:weighted-contract}

All intervals in this section are half-open integer intervals: $[r,s)$ denotes
the points $r,r+1,\ldots,s-1$.  Let $n_0\geq 0$, and let a finite list of
interval isometries generate an equivalence relation on $[0,n_0)$.  Write
$\mathcal C$ for its set of orbits.  The input weight function
$w_0:[0,n_0)\to\mathbb Z^d$, where $d\geq1$, is specified by an ordered
partition into constant runs.  Endpoints and vector entries are binary
integers.  Adjacent runs with equal vectors are coalesced; runs of zero weight
are retained when needed to describe the full domain.

\begin{definition}[Orbit-weight profile]
\label{def:weighted-profile}
For $C\in\mathcal C$, put
\[
 W(C)=\sum_{x\in C}w_0(x).
\]
The compressed profile consists of the pairs $(\mu_v,v)$ for which
\[
 \mu_v=\#\{C\in\mathcal C:W(C)=v\}>0,
 \qquad v\in\mathbb Z^d.
\]
Each distinct vector occurs once, with its multiplicity written in binary.
\end{definition}

Thus components having the same vector remain distinct through their
multiplicity.  A single output record may describe exponentially many
components.  Zero vectors are legitimate, including those produced by
cancellation of positive and negative input weights.  The output does not
select a point or provide a geometric embedding for every component.  For
$n_0=0$, it is the empty profile; the otherwise undetermined vector dimension
can be supplied separately.

The elementary consistency identities are
\begin{equation}
\label{eq:weighted-accounting}
 \sum_v\mu_v=\#\mathcal C,
 \qquad
 \sum_v\mu_v v=\sum_{x=0}^{n_0-1}w_0(x).
\end{equation}
The implementation checks both.  They are useful error detectors, but equal
total mass and orbit count would not by themselves prove the profile correct.
The stronger invariant below supplies that proof.

A simple example shows why multiplicities belong in the contract.  Suppose
one translation, with $1\leq p<n_0$, identifies $x$ with $x+p$ wherever both
points lie in
$[0,n_0)$, and the initial weight is the constant vector $v$.  Write
$n_0=qp+r$, with $0\leq r<p$.  The residue classes give the records
\[
 (p-r,qv),\qquad (r,(q+1)v),
\]
omitting a record of zero multiplicity and combining equal vectors if
necessary.  Thus even an enormous number $p$ of components has at most two
profile records.  If $v=0$, the single record is $(p,0)$; cancellation or a
zero payload never licenses dropping the underlying components.

\subsection{A disjoint-bag invariant}
\label{sec:weighted-invariant}

At an intermediate stage let the current universe be $[0,N)$.  Conceptually
associate to every current point $x$ a bag $B_x$ of original points.  These
bags are pairwise disjoint, and its current vector is
\[
 w(x)=\sum_{y\in B_x}w_0(y).
\]
Empty bags are permitted.  Bags are a proof device and are never stored or
expanded.  For each surviving orbit, the union of its bags is exactly one
original orbit that has not yet been emitted.  Original orbits already emitted
are disjoint from these unions.  Initially, $B_x=\{x\}$.

Identity deletion, reflection trimming, valid periodic merging, and
transmission preserve the equivalence relation on the current point set.
Consequently, they leave all bags and weights unchanged.  This observation
also applies to a valid sharper periodic merger: once its local
partition-preservation statement has been checked, the weight proof needs no
additional assumption about its discovery rule.

The only operations that change weights are the transfer preceding a
truncation and the emission of static points during contraction.  The crucial
transfer rule is the classical one: transfer the \emph{entire carrier range},
even when the subsequent truncation removes only a proper suffix of that
range \cite{AHT2006}.

\subsection{Full-range transfer in half-open coordinates}
\label{sec:weighted-transfer}

Suppose the carrier maps $[a,b)$ isometrically to $[c,N)$, with $a<c$ and
$b-a=N-c$.  The local trace proposes deleting $[L,N)$, where
$c\leq L<N$, after checking that this suffix meets no other pairing.  For
an order-reversing carrier, trimming has made its domain and range disjoint.

First consider disjoint domain and range, so $b\leq c$.  Move each range bag
to its inverse image and clear every bag on $[c,N)$.  For an
order-preserving carrier, put $p=c-a$; the inverse image of a source band
$[r,s)\subseteq[c,N)$ is $[r-p,s-p)$.  For a reflection, whose point map is
$x\mapsto a+N-1-x$, it is
\begin{equation}
\label{eq:weighted-reflection}
 [a+N-s,a+N-r).
\end{equation}
Add the source vector on this image interval.  Because domain and range are
disjoint, these additions do not conflict with clearing the range.

For an overlapping positive translation, $b>c$ and $p=c-a>0$.  Every point
$y\in[c,N)$ has a unique representative
\begin{equation}
\label{eq:weighted-residue}
 \rho(y)=a+((y-a)\bmod p)\in[a,c)
\end{equation}
obtained by repeated partial inverses of the carrier.  These inverses are
defined throughout the descent: intermediate points stay in the range until
the first fundamental block $[a,c)$ is reached.  Transfer the bag at $y$
directly to $\rho(y)$, then clear the entire range.  The algorithm computes
this transfer by division rather than by iterating inverse powers.

Precisely, let a constant source band $[r,s)\subseteq[c,N)$ carry vector $v$.
Define
\[
 q,h=\operatorname{divmod}(s-r,p),\qquad
 u=a+((r-a)\bmod p),\qquad
 t=\min(h,c-u).
\]
Its contribution on the fundamental block is
\begin{equation}
\label{eq:weighted-fold}
 qv\,\mathbf1_{[a,c)}
 +v\,\mathbf1_{[u,u+t)}
 +v\,\mathbf1_{[a,a+h-t)}.
\end{equation}
Empty residual intervals contribute zero.  Indeed, the first $qp$ points
contain exactly $q$ representatives of every residue, and the remaining $h$
points occupy a consecutive cyclic residue interval beginning at $u$.  The
last two terms split that interval at the end of the block.  Formula
\eqref{eq:weighted-fold} uses at most three interval contributions,
independently of the sizes of $q$ and $p$.

\begin{proposition}[Correctness of weighted transfer and emission]
\label{prop:weighted-correctness}
For a valid complete local-reduction trace, the preceding transfers, followed
by static-gap emission, produce the profile in
\cref{def:weighted-profile}.
\end{proposition}

\begin{proof}
An inverse isometry or the map $\rho$ moves each source bag exactly once to
a point in the same current orbit.  Destination bags are replaced by their
disjoint unions with incoming bags.  No original point is duplicated or
discarded.  Clearing the full carrier range therefore leaves every point of
the deleted suffix with an empty bag.  The independently verified unweighted
truncation supplies a bijection between the old and surviving orbits, so the
disjoint-bag invariant is preserved.  Retained points of $[c,L)$ may have empty
bags; the invariant concerns their whole orbit, which still contains the
original mass.

A static gap contains points incident to no current pairing.  Each is an
isolated current orbit, and its bag therefore contains one entire original
orbit.  Intersect the gap with the current constant runs.  An intersection of
length $\ell$ and vector $v$ emits $(\ell,v)$, without visiting its points.
Remove the gap and apply the same increasing coordinate compression to
surviving pairings and runs.  When the validated trace reaches the empty
universe, every original orbit has been emitted exactly once.  Grouping
records with equal vectors gives the claimed profile.
\end{proof}

The implementation performs the interval additions through endpoint
differences.  A contribution $v\mathbf1_{[r,s)}$ places $+v$ at endpoint $r$
and $-v$ at endpoint $s$.  Contributions sharing an endpoint are added
coordinatewise.  Sorting the distinct endpoints and taking prefix sums gives
the constant value between consecutive endpoints.  Neighboring output bands
with equal values are immediately coalesced, including zero-valued bands.
The old weight below $c$ is included in the same sweep, whereas old range
weights are used only as transfer sources.  This simultaneous construction
avoids accidentally transferring a value that was just added during the
current operation.  In particular, periodic folding is a push-forward of
the old weights, not repeated in-place accumulation along a carrier chain.

\subsection{The additive bound on weight breakpoints}
\label{sec:weighted-breakpoints}

It is insufficient merely to note that each source run produces at most three
contributions: such a statement alone would permit a multiplicative increase
at every step.  The polynomial bound depends on the disappearance and
relocation of old breakpoints, as in AHT Lemma~15 \cite{AHT2006}.

\begin{lemma}[Whole-range breakpoint bound]
\label{lem:weighted-bands}
One full-range transfer increases the number of maximal constant runs by at
most four.  Truncation and contraction do not increase it.  Consequently,
after $s$ transfers, the number of active runs is at most $m+4s$, where $m$
is the initial number.
\end{lemma}

\begin{proof}
For $N>0$, write
\[
 D(w)=\{j:1\leq j<N,\ w(j)\ne w(j-1)\}.
\]
The run count is $1+|D(w)|$.  Split old discontinuities into those strictly
below $c$ and those strictly above $c$, leaving a possible discontinuity at
$c$ aside.  Discontinuities below $c$ remain potential discontinuities of the
unchanged old weight.  All those strictly inside the range disappear there,
since that entire range becomes zero.

In the disjoint case, a range discontinuity at $j$ has at most one image:
$j-p$ for a translation or $a+N-j$ for a reflection.  Away from these images,
inherited discontinuities below $c$, and the endpoints $a,b,c$, both the old
weight and the added weight are locally constant.  Thus there are at most
four exceptional boundary locations in addition to inherited or relocated
old discontinuities.

For completeness, the modular case can be checked directly by finite
differences.  Extend the range weight by zero to all integers, writing
$z(j)=w(j)$ on $[c,N)$ and $z(j)=0$ otherwise.  The incoming periodic weight is
\[
 H(x)=\sum_{\ell\in\mathbb Z}z(x+\ell p),\qquad a\leq x<c.
\]
Each sum is finite.  For an interior block position,
\[
 H(x)-H(x-1)
 =\sum_{\ell\in\mathbb Z}
   \bigl(z(x+\ell p)-z(x-1+\ell p)\bigr).
\]
The nonzero differences of $z$ occur only at old interior range
discontinuities and at $c,N$.  Each interior discontinuity has only one
residue image in $[a,c)$; the boundary $c$ maps to $a$.  Apart from inherited
and relocated old discontinuities, possible new boundaries are contained in
$\{a,c,\rho(N)\}$, using the same residue formula at the boundary $N$.
Again there are at most four.  Vector addition or signed
cancellation can remove discontinuities but cannot create any outside this
set.  Counting old discontinuities only once proves the asserted bound.

Restricting to a shorter prefix cannot create a new run.  For contraction,
delete the gap pieces from the ordered sequence of runs and concatenate the
surviving pieces.  Pieces of any one old run become adjacent if separated only
by deleted points, and equal neighboring values can be coalesced.  The number
of runs therefore cannot increase.  Induction gives $m+4s$.
\end{proof}

This argument explains why transferring only $[L,N)$ is an inadequate
implementation shortcut for the stated complexity proof.  When $L>c$,
breakpoints in the remaining range could survive while translated or residue
images of other breakpoints are added.  The one-for-one relocation count
would no longer apply.  Such a shortcut may conserve total weight, but that
fact does not establish a polynomial bound on its representation.  The
implementation instead performs full-range transfer and checks the
at-most-four increase before deleting the zero suffix.

\subsection{Trace length, integer size, and computational cost}
\label{sec:weighted-complexity}

Let $P$ be the number of operations in a supplied complete trace, $k$ the
initial number of pairings, and
\[
 B=\max_{x,j}|w_0(x)_j|,\qquad M=m+4P,
\]
with $B=0$ for the empty input.
The number of pairings never increases.  By
\cref{lem:weighted-bands}, at most $M$ weight runs are active.  One transfer
creates $O(M)$ interval contributions and combines them through a sorted
endpoint-difference sweep.  There is no loop over an interval length, a
periodic quotient, or a component multiplicity.

There are at most $2k+1$ static gaps.  Overlaying them with the run partition
during a contraction produces $O(M+k)$ emission records.  Hence
$R=P(M+2k+1)$ is a conservative bound on their total number before grouping.
Including the elementary independent trace checker, coordinate updates, and
sorting equal profiles, a coarse arithmetic-operation bound under ordinary
expected dictionary-operation costs is
\begin{equation}
\label{eq:weighted-arithmetic-cost}
 O\!\left(
 dm+k+P\bigl((k+1)^2+dM\log(M+1)\bigr)
       +dR\log(R+1)
 \right).
\end{equation}
The implementation uses Python dictionaries for endpoint events and a
counter for profile grouping.  Arbitrary large integers can have colliding
hashes, so the displayed bound is not a deterministic worst-case bound for
those containers.  With elementary quadratic bounds for their assembly and
grouping, a deterministic bound for a nonempty completed trace is
\[
 O\!\left(dm+k+P(k+1)^2+dR^2\right).
\]
Indeed $R=P(M+2k+1)$ with $P\geq1$ dominates both the per-transfer event-map
work and final grouping; the empty trace is trivial.  Both bounds are
polynomial.  No optimized exponent for weighted AHT is asserted.  Their
purpose is to make the dependence on compressed input and explicit trace
length transparent without assuming adversarial hashing is free.

Signed coordinates also have controlled bit length.  From the disjoint-bag
invariant, for every current point and coordinate,
\begin{equation}
\label{eq:weighted-bit-bound}
 |w(x)_j|
 \leq\sum_{y\in B_x}|w_0(y)_j|
 \leq\sum_{y<n_0}|w_0(y)_j|
 \leq n_0B.
\end{equation}
The same absolute-mass argument bounds endpoint-difference intermediates up
to an absolute constant factor.  Their bit lengths are therefore
$O(\log(n_0+1)+\log(B+1))$.  Endpoints, valid transmission powers, quotients,
and multiplicities have $O(\log(n_0+1))$ bits.  Substituting standard integer
arithmetic costs into \eqref{eq:weighted-arithmetic-cost} gives polynomial bit
complexity in $P,k,m,d,\log(n_0+1),\log(B+1)$.

More explicitly, disjointness gives the stronger coordinatewise bound
\[
 \sum_{x<N}|w(x)_j|
 \leq \sum_{y<n_0}|w_0(y)_j|.
\]
For a source band of length $\ell$, the whole-period coefficient in
\eqref{eq:weighted-fold} is at most $\ell$, and there are at most two residual
contributions.  The total absolute coefficients used to assemble endpoint
events are therefore at most a fixed constant times the source's absolute
mass.  Including retained old bands preserves such a constant-factor bound.
This justifies the intermediate-size assertion even before cancelling
endpoint contributions have been combined; signed cancellation is never
needed to keep the arithmetic representation polynomial.

The distinction between trace checking and trace discovery matters.  The
classical AHT scheduling analysis supplies a polynomial bound on $P$ in
$k$ and $\log(n_0+1)$, yielding the weighted polynomial-time conclusion of
AHT Theorem~16 \cite{AHT2006}.  For another scheduler, weighted replay inherits
whatever trace-length bound has been established for that scheduler.  Valid
local reductions alone do not prove that an arbitrary strategy finds a short
trace.  Conversely, once a short valid trace is available, its weighted
evaluation does not depend on how it was discovered.

\subsection{Certificate binding and interruption semantics}
\label{sec:weighted-certificates}

The replay entry point first invokes the independent unweighted checker.
This binds the trace to the supplied universe and normalized ordered pairing
list, and checks every local operation and the terminal orbit count.  The
weighted replay then follows those validated operations without rerunning
the producer's scheduler.  It accepts both supported trace versions, whose
different periodic-merger conditions are checked by that verifier.

The optional returned certificate is this unweighted trace.  It does not
itself bind a new weight function: the same trace can correctly be reused for
many different weights.  A reproducible weighted assertion must therefore
supply the interval system, full run partition, dimension, claimed profile,
and trace together.  A consumer reconstructs or validates the weights and
replays them before comparing the claimed profile.  Geometric conclusions
require a separate justification that these input weights measure the desired
geometric quantities.

The public wrapper can produce a trace or reuse an existing one.  Its
\texttt{max\_cycles} parameter bounds producer cycles, not bit operations or
elapsed time.  An exhausted producer returns an incomplete result with no
orbit count, profile, or certificate.  A reused trace has no asserted producer
cycle count; its evaluation is controlled through the cooperative
\texttt{check} callback instead.  The callback is used throughout production,
validation, and replay, and its exceptions propagate.  Weight or trace
rejections use the dedicated \texttt{WeightedOrbitError} exception, allowing
a caller to distinguish invalid evidence from cancellation, including a
callback that raises an ordinary \texttt{ValueError}.

These interfaces support exact component statistics while preserving the
logical distinction between a completed proof, rejected evidence, and an
interrupted computation.  They supply a polynomial component-analysis layer;
the search or hierarchy needed to obtain a complete unknot-recognition bound
remains a separate mathematical obligation.
